\emph{Step 6. The equation.} By Step~5, $\int\nu^*(\theta-W)=0$ with integrand $\ge0$. It follows that $\nu^*(\theta-W)=0$ almost
everywhere and, since $\nu^*>0$, that $W=\theta$. That is, $\nu^*=e^{-\theta}\,\Pi[\nu^*]$.

\emph{Step 7. The bounds.} By the equation and Lemma~\ref{lem:fubini}, for almost every $r\in I_R$, \[
 m^*_R(r)=\lambda\,u^*_R(r)\int_0^1u^*_C(c)\,u^*_D(r-c)\,u^*_A(r+c)\,dc\ \le\ \lambda\bigl(1-m^*_R(r)\bigr),
\] which gives $m^*_R\le\lambda/(1+\lambda)$. The same argument works for the other families, since the integral runs over a fiber of
parameter length $\le1$. Each factor in $\Pi[\nu^*]$ therefore lies between $1/(1+\lambda)$ and $1$, giving
$\lambda(1+\lambda)^{-4}\le\nu^*\le\lambda$. Integrating the equation gives $\rho=\lambda\int\Pi[\nu^*]$.

\emph{Step 8. A supergradient.} Let $\rho'\in(0,1)$, $\nu'\in\cA_{\rho'}$ with $s[\nu']>-\infty$, and let
$\varphi(t)=s[\nu^*+t(\nu'-\nu^*)]$, which is concave on $[0,1]$. By Step~7, $\log\nu^*$ and the $\log u^*_f$ are bounded, so the
calculation in Step~2 applies with all terms finite. The masses now differ, but the constants again cancel. Their contributions are
$-(\rho'-\rho)$ from the entropy, $+4(\rho'-\rho)$ from the four families, and $-3(\rho'-\rho)$ from the term $-3\rho$. This leaves
$\varphi'(0^+)=\int(\nu'-\nu^*)W=\theta(\rho'-\rho)$, and $s[\nu']=\varphi(1)\le\varphi(0)+\varphi'(0^+)$. Taking the supremum over
$\nu'$ gives the inequality in~(d). A function that lies below a line touching it at each point is concave. If $\rho_1<\rho_2$, adding
the inequalities at $\rho_1$ and $\rho_2$ gives $(\theta(\rho_1)-\theta(\rho_2))(\rho_2-\rho_1)\ge0$.

\emph{Step 9. The derivative.} If $|\sigma-1|$ is small, $\sigma\nu^*\in\cA_{\sigma\rho}$ because its projections are
$\le\sigma\lambda/(1+\lambda)<1$. The function \[
 \psi(\sigma)=s[\sigma\nu^*]=-3\sigma\rho-\sigma\int\nu^*\log\nu^*-\sigma\rho\log\sigma+\sum_f\int_{I_f}G(\sigma m^*_f)
\] is of class $C^2$ near $\sigma=1$, with $\psi'(1)=\int\nu^*W=\rho\theta$ and
$|\psi''(\sigma)|\le\rho/\sigma+\sum_f\int(m^*_f)^2/(1-\sigma m^*_f)$, which is bounded. Thus
$\sopt(\rho+h)\ge\psi(1+h/\rho)\ge\sopt(\rho)+\theta h-Ch^2$ for small $|h|$. Together with Step~8 this gives
$\sopt(\rho+h)=\sopt(\rho)+\theta h+O(h^2)$, so $\sopt$ is differentiable with $\sopt{}'=\theta$. The derivative of a differentiable
concave function on an open interval is continuous. \end{proof}

\begin{proposicion}\label{prop:concava} The function $\sopt$ is strictly concave on $(0,1)$. Consequently, $\theta=\sopt{}'$ is
continuous and strictly decreasing, and $\lambda=e^{-\theta}$ is continuous and strictly increasing. Moreover,
$\lambda(1+\lambda)^{-4}\le\rho\le\lambda/(1+\lambda)$. It follows that $\lambda(\rho)\to0$ as $\rho\to0$ and
$\lambda(\rho)\ge\rho/(1-\rho)\to\infty$ as $\rho\to1$. \end{proposicion}

\begin{proof} The functional $s[\nu]=-3\int\nu-\int\nu\log\nu+\sum_f\int G(m_f)$ is concave on the convex union of the classes
$\cA_\rho$, and its entropy term is strictly concave. If $\sopt$ were affine on $[\rho_1,\rho_2]$, the mixture
$\bar\nu=\frac12(\nu_{\rho_1}+\nu_{\rho_2})\in\cA_{(\rho_1+\rho_2)/2}$ would satisfy
$s[\bar\nu]>\frac12(\sopt(\rho_1)+\sopt(\rho_2))=\sopt(\frac{\rho_1+\rho_2}2)$, which is impossible. The derivative of a differentiable
strictly concave function is strictly decreasing, and it is continuous by Theorem~\ref{thm:max}(d). Integrating the bounds on
$\nu_\rho$ and $m_R$ from Theorem~\ref{thm:max}(c) over the square gives $\lambda(1+\lambda)^{-4}\le\rho\le\lambda/(1+\lambda)$. By
monotonicity, $\lambda$ has a limit as $\rho\to0$, and the first inequality forces that limit to be zero. The second gives
$\lambda\ge\rho/(1-\rho)$. \end{proof}

Theorem~\ref{thm:max} and Proposition~\ref{prop:concava} prove Theorem~\ref{thm:euler} for $\rho<1$. The case $\rho=1$ is treated in
Proposition~\ref{prop:masa_uno}.

\begin{lema}[Representation of the vacancies]\label{lem:vacancias} Let $0<\rho<1$, $\nu=\nu_\rho$, and $\lambda=\lambda(\rho)$. For
each family $f$ and each $t\in I_f$, let $J_f(t)$ be the integral over the fiber of $f$ over $t$, with respect to the parameter
measure, of the product of the vacancies of the other three families. For example, $J_R(r)=\int_0^1u_C(c)\,u_D(r-c)\,u_A(r+c)\,dc$.
Then $0\le J_f\le1$, $J_f$ is continuous on $I_f$, and \[
 u_f=\frac1{1+\lambda J_f},\qquad m_f=\frac{\lambda J_f}{1+\lambda J_f}\qquad\text{almost everywhere}.
\] We use these continuous representatives in what follows. In particular, $u_f\ge1/(1+\lambda)$, $u_D(\pm1)=1$, and $u_A(0)=u_A(2)=1$. \end{lema}

\begin{proof} The integral $J_f(t)$ is defined for every $t$. On a fiber of $f$, the coordinate of each other family is an affine
function of the parameter with slope of absolute value one or two, so the composition is independent of the choice of representatives.
By the Euler equation and Lemma~\ref{lem:fubini}, $m_f=\lambda u_fJ_f$ almost everywhere, as in Step~7, and the formula follows by
rearranging. The factors lie in $[0,1]$, and the fibers have length at most one. When passing from $t$ to $t'$, the fiber changes on a
set of length $\le2|t-t'|$, and the arguments of the other factors are shifted by $|t-t'|$ (on diagonals and antidiagonals, after a
change of variable with Jacobian $\frac12$). With the $m_g$ extended by zero, \[
 |J_f(t)-J_f(t')|\le2|t-t'|+2\sum_{g\ne f}\ \sup_{|h|\le|t-t'|}\|m_g(\cdot+h)-m_g\|_{L^1(\mathbb R)},
\] which tends to zero by continuity of translations in $L^1$. At the endpoints of $I_D$ and $I_A$, the fiber is a point and $J_f=0$. \end{proof}

\begin{proposicion}[Continuity with respect to mass]\label{prop:continuidad} \begin{enumerate}[(a)] \item The map $\rho\mapsto\nu_\rho$
is continuous from $(0,1)$ to $L^1([0,1]^2)$, and $\rho\mapsto m_{f,\rho}$ is continuous from $(0,1)$ to $L^2(I_f)$ for all four
families. \item For each compact interval $[\rho_1,\rho_2]\subset(0,1)$, the family $\{m_{f,\rho}:\rho\in[\rho_1,\rho_2]\}$ of
continuous representatives is equicontinuous. Hence $(\rho,t)\mapsto m_{f,\rho}(t)$ is uniformly continuous on $[\rho_1,\rho_2]\times
I_f$. \end{enumerate} \end{proposicion}

\begin{proof} (a) Let $\rho_j\to\rho$, $\nu_j=\nu_{\rho_j}$, and $\nu=\nu_\rho$. Since $\sopt$ is continuous and the line terms are
bounded, the entropies of the $\nu_j$ are bounded, and every subsequence has a further subsequence that converges weakly in $L^1$ to
some $\bar\nu$ (Step~1). By Lemma~\ref{lem:fubini}, $\bar\nu\in\cA_\rho$, since its mass is $\rho$ and $\int m_f[\bar\nu]\,g=\lim\int
m_f[\nu_j]\,g\le\int g$ for every bounded $g\ge0$. Each term of $s$ (the mass term, the entropy, and the four line terms) is concave
and upper semicontinuous in $L^1$ and hence weakly upper semicontinuous by Mazur's lemma. Thus $s[\bar\nu]\ge\limsup
s[\nu_j]=\sopt(\rho)$, and $\bar\nu=\nu$ by uniqueness. We therefore have $\nu_j\to\nu$ weakly. Since the sum converges and the limit
superior of each term is at most its value at $\nu$, each term converges. With $H(z)=z\log z$ and $H''=1/z\ge1/(u+v)$ between $u$ and $v$,
\[
 \frac{H(u)+H(v)}2-H\Bigl(\frac{u+v}2\Bigr)\ge\frac{(u-v)^2}{8(u+v)}\qquad(u,v\ge0).
\] Integrating with $u=\nu_j$, $v=\nu$ and using weak lower semicontinuity of $\int H$ in $\frac{\nu_j+\nu}2\to\nu$ gives
$\int\frac{(\nu_j-\nu)^2}{\nu_j+\nu}\to0$, and Cauchy--Schwarz gives
$\|\nu_j-\nu\|_1^2\le(\rho_j+\rho)\int\frac{(\nu_j-\nu)^2}{\nu_j+\nu}\to0$. For the line terms, $G''\le-1$ on $[0,1)$ gives $\int
G(\frac{m_j+m}2)\ge\frac12\int G(m_j)+\frac12\int G(m)+\frac18\|m_j-m\|_2^2$, and upper semicontinuity in $\frac{m_j+m}2\to m$ gives
$\|m_j-m\|_2\to0$.

(b) By Lemma~\ref{lem:vacancias}, $m_{f,\rho}=\lambda J_f/(1+\lambda J_f)$, with the continuity bound for $J_f$ from its proof. By~(a),
$\{m_{g,\rho}:\rho\in[\rho_1,\rho_2]\}$ is compact in $L^1$ and therefore has a uniform modulus of continuity under translation
(Fréchet--Kolmogorov~\cite{Brezis}). Since $\lambda(\rho)\le\lambda(\rho_2)$ and $z\mapsto\lambda z/(1+\lambda z)$ is
$\lambda$-Lipschitz, the family is equicontinuous. Equicontinuity, boundedness, and convergence in $L^2$ give uniform convergence by
Arzelà--Ascoli, and joint continuity follows. \end{proof}
